\chapter*{Notations}

\begin{itemize}
    \item $\mathcal{X}$ désigne l'ensemble des entrées (observations). $\mathcal{X}^m$ est l'ensemble des $m$-uplets $(x_1, \dots, x_m)$
    dont chaque terme est un élément de $\mathcal{X}$, on pose $\mathcal{X}^{*} = \bigcup_{m=1}^{+\infty} \mathcal{X}^{m}$.
    En d'autres tertmes, 
    \[
    \mathcal{X}^m = 
    \begin{bmatrix}
        x_1, \dots, x_m
    \end{bmatrix}
    \quad x_i \in \mathcal{X}, \forall i \in \{1;m\}
    \]
    \[
    \quad \mathcal{X}^{*} =\bigcup_{m=1}^{+\infty} \mathcal{X}^{m} =
    \begin{bmatrix}
        x_{1,1} & \dots & x_{1,m} \\
        \vdots   & \dots & \vdots \\
        x_{n,1} & \dots & x_{n,m}
    \end{bmatrix}
    \]
    \newline
    \item $|\mathcal{X}|$ désigne le cardinal de $\mathcal{X}$ \newline 
    \item $\mathcal{Y}$ désigne l'espace des sorties. On a aussi $\mathcal{Y}^m$ est l'ensemble des $m$-uplets $(y_1, \dots, y_m)$
    dont chaque terme est un élément de $\mathcal{Y}$, on pose $\mathcal{Y}^{*} = \bigcup_{m=1}^{+\infty} \mathcal{Y}^{m}$. 
    Il vient :
    \begin{equation}
        S = ((x_i, y_i)_{i \in \{1;m\}}) \in (\mathcal{X} \times \mathcal{Y})^{*} \nonumber
    \end{equation} \newline
    \item $L$ désigne la fonction de coût \newline
    \item $m(\mathcal{X}, \mathcal{Y})$ est l'ensemble des fonctions mesurables \footnote{
        ($\mathcal{X}, m_{\mathcal{X}}$) et ($\mathcal{Y}, m_{\mathcal{Y}}$) sont des espaces mesurables
    } \newline
    \item $m_{\mathcal{H}}(\epsilon, \delta)$ est la fonction de complexité d'échantilon. \newline
    \item $(ERM)_S$ est l'hypothèse qui minimise le risque empirique sur $S$.
    \item ($\Omega, \mathcal{E}, \mathbb{P}$) désigne un espace probabilisé \newline
    \item $\mathbf{A}$ désigne une matrice. \newline
    \item $\mathcal{M}_{n,p}(\mathbb{K})$ désigne l'espace des matrices de $n$ lignes, $p$ colonnes à coefficients dans $\mathbb{R}$
\end{itemize}

\chapter*{Rappels}

\noindent

\chapter{Modèle formel de l'apprentissage statistique supervisé}
\section{Définitions}

\begin{itemize}

\item \noindent Soient $\mathcal{X}, \mathcal{Y}$ donnés. La fonction de coût $L$ est définie par : 

\[
\begin{aligned}
    L : \mathcal{Y} \times \mathcal{Y} &\to \mathbb{R}^{+} \\
    (y,y') &\mapsto L(y,y')
\end{aligned}
\]

\noindent C'est une fonction arbitraire. En général, on prend $L(y,y') = |y - y'|^2$ (méthode OLS) avec $y' = \hat{y} = h(x)$, 
$h$ l'hypothèse que l'on fait sur $x \in \mathcal{X}$ pour estimer la valeur $y$ i.e. $\hat{y}$ est un estimateur de $y$. $L$ est donc le coût
de décision de $h$ sur $x$. \newline

\item $h \in \mathcal{H} \subset m(\mathcal{X}, \mathcal{Y})$ avec $\mathcal{H}$ l'ensemble des hypothèses faites sur $x$ pour estimer $y$. \newline

\item $S = ((x,y_i))_{1 \leq i \leq m} = ((x_1,y_1), \dots, (x_m,y_m)) \in (\mathcal{X} \times \mathcal{Y})^{*}$  est une séquence d'objets labellisés.
Les $x_i$ sont les données observées et les $y_i$ les valeurs à estimées. $S$ représente le jeu de données d'entraînement \newline

\item Un Algorithme Automatique Supervisé (AAS) est une fonction tel que : 

\[
\begin{aligned}
A : (\mathcal{X} \times \mathcal{Y})^{*} &\to \mathcal{H} \\
S &\mapsto A(S) = h_s
\end{aligned}
\]

\item Risque réel :

\begin{align}
    \mathcal{R}(h) &= \mathbb{E}[L(y,y')] \\
                   &= \int_{\mathcal{X} \times \mathcal{Y}} L(y, h(x))p(x,y)dxdy \nonumber\\
                   &= \int_{\mathcal{X} \times \mathcal{Y}} L(y, h(x))p(y | x)p(x)dxdy \nonumber
\end{align}


\item Risque empirique :

\begin{equation}
    \mathcal{R}_s(h) = \hat{\mathcal{R}}(h) = \frac{1}{m} \sum_{i=1}^{m} L(h(x_i), y_i)
\end{equation}


\newpage

\noindent C'est le seul risque dont on a accès (quantité d'information disponible limitée et finie). On cherche à minimiser 
le risque réel $\mathcal{R}$ à l'aide de la décision prise parmis $\mathcal{H}, h_s = A(S)$, à savoir :

\begin{equation}
    \mathcal{R}(h_s) = \min_{h \in \mathcal{H}}\mathcal{R}(h) \Longleftrightarrow 
    (\text{ERM})_{\mathcal{H}}  = h_s \in \text{arg}\min_{h \in \mathcal{H}} \mathcal{R}_s(h)
\end{equation}

\noindent \textbf{Propriété :} $h \in \mathcal{H}, \mathcal{R}_s(h) \xrightarrow[m \to \infty]{p.s.} \mathcal{R}(h)$ et en moyenne quadratique. 
\footnote{Attention : $A(S)$ est une variable aléatoire donc $\mathcal{R}(A(S)) = \mathcal{R}(h_s)$ est une VA fonction de $S$, on ne peut pas appliquer la loi des grands nombres.}

\end{itemize}

\section{Capacité d'apprentissage PAC}

\noindent Soient $\mathcal{X}, \mathcal{Y}, L, \mathcal{H}$ fixés. $\mathcal{H}$ a la capacité PAC 
 \footnote{Probablement Approximativement Correct. Probablement car la probabilité de l'évènement est non négligeable ($\geq 1 - \delta$),
 Approximativement car le minimum global n'est pas atteint mais il s'en rapproche d'une distance au plus $\epsilon$ ($\geq \epsilon$)}
quand $\exists A, \exists m_{\mathcal{H}} : ]0;1[^{2} \to \mathbb{N}, \forall P, \forall (\epsilon, \delta) \in  ]0;1[^{2}, 
\forall m \ge m_{\mathcal{H}}(\epsilon, \delta), \forall S$

\begin{equation}
    P(\mathcal{R}(A(S)) - \min_{h \in \mathcal{H}}\mathcal{R}(h) \leq \epsilon) \geq 1 - \delta 
    \Longleftrightarrow P(\mathcal{R}(h_S) - \min_{h \in \mathcal{H}}\mathcal{R}(h) \leq \epsilon) \geq 1 - \delta \quad (*) \nonumber
\end{equation}

\noindent De même :

\begin{equation}
    m_{\mathcal{H}}(\epsilon, \delta) = \min \{m \in \mathbb{N} \quad \text{tel que} \quad \forall P, \forall S, \quad \text{(*) est vérifiée}\}
\end{equation}

\noindent $m_{\mathcal{H}}(\epsilon, \delta)$ est la taille minimale du \textit{dataset} qui permet à l'algorithme d'apprendre correctement. C'est
la fonction de complexité de $\mathcal{H}$ en mémoire ou en échantilon (données). Plus $\mathcal{H}$ est complexe, plus il faut de données. De même, plus $\epsilon, \delta$ sont petits,
plus il faut de données. Cela veut dire que l'algorithme retourne un bon modè-le avec une grande probabilité. En général, pour une classe de VC-dimension $d$, on prend 
$m_{\mathcal{H}}(\epsilon, \delta) \approx \frac{d + log(\frac{1}{\delta})}{\epsilon}$\newline

\noindent \textbf{Propriété :} Soit $\epsilon > 0$, $S$ est \textbf{$\epsilon$-représentatif} si et seulement si pour $h_s = (ERM)_{\mathcal{H}}(S)$,

\begin{equation}
    \mathcal{R}(h_s) - \min_{h \in \mathcal{H}} \mathcal{R}(h) \leq \epsilon
\end{equation}

\section{Convergence uniforme}

\noindent $\mathcal{H}$ a la propriété de convergence uniforme quand $\exists m_{\mathcal{H}}^{CU} \in ]0;1[^{2} \rightarrow \mathbb{N},
\forall P, \forall (\epsilon, \delta) \in  ]0;1[^{2}, \forall m \geq m_{\mathcal{H}}^{CU}(\epsilon, \delta), \forall S,$

\begin{equation}
    P(\text{S est espilon-représentatif}) \geq 1 - \delta \Longleftrightarrow
    P(\sup_{h \in \mathcal{H}} | \mathcal{R}(h) - \mathcal{R}_S(h) | \leq \epsilon) \geq 1 - \delta \quad (**) \nonumber
\end{equation}

\noindent \textbf{Propriété :} Si $\mathcal{H}$ a la propriété de CU avec $m_{H_{CU}}^{CU}$ alors $\mathcal{H}$ a la capacité PAC par $A = (ERM)_{\mathcal{H}}$ et 

\begin{equation}
    m_{H}(\epsilon, \delta) \leq m_{H}^{CU}(\frac{\epsilon}{2}, \delta)
\end{equation}

\newpage

\section{Capacité d'apprentissage PAC pour des ensemble finis}

\noindent \textbf{Lemme d'Hoeffding (1962)} : pour $(X_i)_{1 \geq i \geq m}$ VA i.i.d., $\mathbb{E}[X_i] = \mu, \exists a < b, a,b \in, \mathbb{R}$
tel que $P(X_i \in [a;b]) = 1$, alors, $\forall \epsilon > 0$,

\begin{equation}
    P(| \frac{1}{m}\sum_{i=1}^{m}X_i - \mu|) \leq 2e^{ \frac{-2m\epsilon^2}{(b-a)^2}}
\end{equation}

\noindent On appelle cette inégalité l'inégalité de concentration.